% !TeX root = ../papers/nc-p3.tex
\section{有限点窓の三文字評価の階数差と点列のイデアル}
\label{sec:finite-three-slot-evaluations}
\begin{theorem}\label{thm:finite-three-slot-evaluations}
標準有限Frobenius候補の四辺が有限とする。共通点環 $D$ における連続三位置の座標を
$U,V,W$ とし、単元倍の違いを許す。反対位置 $U,W$ が射影同定されれば
$\dim UVW=16$ である。同定されないときは19または20であり、
階数19での三積の唯一の零化汎関数の支持スキームの長さは9以上である。
従って長さ3の経路から点列への評価の核は、順に次元4、1、0で尽きる。
\end{theorem}
\begin{proof}
$W=U$ に規格化できる場合は辺の六式の完全交叉のKoszul分解を $(2,1)$ でひねる。
二つの斉次次数がともに非負なのは初めの二項のみで、最後まで高次コホモロジーから
大域切断の核へ入る項はない（両次数とも $-4$ 以下となるのは第六項だけで全次数0）。
従って像の次元は $10\cdot4-6\cdot4=16$ である。

他の場合は定理\ref{thm:finite-opposite-product-gap}から $\dim UW=16$ である。
$D$ の完全積対合の跡を一つ選び、$E=(UW)^\perp$ とすると $\dim E=4$。
元 $a\ne0$ の支持は $Y_a=\operatorname{Spec}(D/\operatorname{ann}(a))$ とし、
その長さを $m_a$ とする。各座標空間は単元を含み、任意の長さ $m$ の部分スキームへの
制限次元には、$m\ge1,2,4,8$ に応じてそれぞれ $1,2,3,4$ の下界がある。
これは各位置で点・直線・平面の切断長がそれぞれ1,3,7以下であるためである。
$a\in E$ なら $Y_a$ 上の $U,W$ が $a$ による完全跡で直交する。
上の下界を足すと $m_a\ge6$ が必要である。

三積の零化元を $g$ とすると条件は $gV\subset E$ である。$0<m_g\le7$ と仮定し、
$r=\dim(V|_{Y_g})$ とする。$Y_g$ の長さ $r-1$ の部分スキーム上で零で、$Y_g$ 上では非零な
$v\in V$ が存在する（Artin環のイデアルの組成列を使う）。すると $gv\ne0$、
$m_{gv}\le m_g-r+1\le5$ となり $E$ の下界に反する。$m_g=8$ も不可能である。
そのとき $Y_g$ 上の $UV$ の次元は補題\ref{lem:finite-adjacent-strict-half}で5以上、
$W$ は4であり直交できない。従って $m_g\ge9$、特に $gV=E$ となる。

さらに $g'\ne0$ も三積を消せば、両者が生成する $D$ のイデアルは
$DE=(g)=(g')$ である。ゆえに $Y=Y_g$ 上で $g'=ug$ かつ $u$ は単元で、
$u(V|_Y)=V|_Y$。このような $u$ は定数しかない。実際非定数なら剰余値の一つを引き、
$V|_Y$ を保つ非零・非単元 $a$ を取る。$r=\dim a(V|_Y)$ は、その支持への $V$ の制限次元に等しい。
また $Y$ 上の $a=0$ への制限次元は高々 $4-r$ である。
支持と零点スキームの長さの和は $\operatorname{length}Y$。どちらも空でないため
$r=1,2,3$ であり、先の部分空間切断長の上界からその和は
$1+7,3+3,7+1$ 以下である。これは $m_g\ge9$ に反する。
従って零化空間は高々一次元である。最後に候補の三次重なりは次元4なので
経路空間の次元は $4^3-2\cdot 4\cdot6+4=20$ である。
\end{proof}

\begin{corollary}\label{cor:finite-asymmetric-point-ideal}
両方の反対対で射影同定がない有限窓とし、$A$ から共通点列の切断の積への評価の核を $I$ とする。
$I$ の各始点での右イデアルは次数3,4の部分で生成され、左についても同様である。
各始点で三次生成元は高々一本である。全次数のAS性を仮定しない。
\end{corollary}
\begin{proof}
始点を固定した像を $S_n\subset D$ と書く。定理\ref{thm:finite-four-slot-rank}により $S_4=D$、
以後も各座標に単元があるので $S_n=D$ である。$I$ の長さ $n>1$ の右極小生成元の空間は、
$A$ が二次表示なので、点列の切断に関する次の三項の中央ホモロジーである。
\[
 S_{n-2}\otimes R_j\longrightarrow S_{n-1}\otimes V_{j+1}\longrightarrow S_n ,
 \qquad j=i+n-2 .
\]
実際これは $S$ と終点単純の左自由表示をテンソルした第一Torである。
全項の係数空間が $D$ なら点グラフ上の関係の左評価は階数3の直和因子であり、
次の点での評価の核を張るので完全である。これで $n\ge6$ は済む。

$n=5$ で $S_3$ が超平面の場合だけ残る。完全跡で $S_3^\perp=kg$ とし、$Y=Y_g$ とする。
$D$ 上の全体の写像 $\alpha:D\otimes R_j\to D\otimes V_{j+1}$ の像に、
$S_3\otimes R_j$ の像が届かないとすると、両者の余核の双対を比較して、
$0\ne g\otimes c\in\operatorname{im}\alpha^*$、$c\in R_j^*$ が存在する。
実際この制限で消える汎関数の全体は $kg\otimes R_j^*$ である。
相補切断により $R_j^*\simeq R_{j+2}$ とする。$\alpha^*$ の像は
$w(x_j,-,-,-)$ の次文字による縮約で張られるので、どの元も最後の点 $x_{j+3}$
で評価すると零である。これは四移動が点スキーム上で恒等で、辺 $x_{j+3}\to x_j$ が成立するからである。
従って $c$ の末尾支持の文字は $Y$ 上で全て零になる。ところが前定理より
$Y$ の長さは9以上で、一平面に入らないので全て零、矛盾となる。逆向きも同じである。
\end{proof}

\begin{corollary}\label{cor:finite-full-triple-regular-presentation}
十五次元混合成分の幾何学的一般点では全ての三文字評価が同型である。
より一般に全三文字評価が同型の有限窓では、各始点の点列イデアルは四次生成であり、
その極小生成元の間の右関係は全て五次で生成される。左でも同様である。
この関係の本数は各位置で高々24。標準正則なら四次元の窓ごとに
生成元は15個、関係は24個、これらの間の第二関係は六次の10個で尽き、それより高い関係はない。
\end{corollary}
\begin{proof}
定理\ref{thm:marked-twenty-rational-invariants}で構成した二十印の持上げにおいて、
各位置の三文字の全単項式の評価行列を取る。付属検算は素数1009の整数環を十四桁で切った
環上で単元ピボットを消去する。全四位置でピボットは20あり、二位置では全て単元、
残る二位置でも付値は16個が0で4個が1である。従って行列式小行列は標数零の一般側で非零となる。
階数20は開条件なので同成分の一般点に成立する。

各 $S_n$ は $n\ge3$ で $D$ である。三次核が零なので先の系で四次生成である。
終点単純の左Koszul候補複体 $K$ と $S$ をテンソルする。$K$ は二次表示と全三次重なりを使うため、
$\operatorname{Tor}_2(S,k)$ はこの第二ホモロジーの商である。実際、真の第三自由項には
全ての三次重なりが入り、それ以外の生成元があれば境界が増えるだけである。
$n\ge6$ では項 $1,2,3$ の係数が全て $D$ なので、点列での階数 $1,3,3,1$ から
第二ホモロジーは零である。四次生成の極小関係の次数は5以上なので五次だけが残る。

$n=5$ では $S\otimes K$ の項の次元が右端から $20,80,120,40,4$。
最後の写像は単射で、第三ホモロジーも零であることを示す。
全係数を $D$ に広げると第三項の微分の核は、前の点の四座標を並べた列の $D$ 倍である。
これは点列での分裂した完全性と、第四微分がその座標列であることによる。
元の第三項では各係数が $UV$（初めの二位置）に入るので、核の倍率 $a$ の条件は $aV\subset UV$。
$(2,2)$ の完全跡に関し $UV=(UV)^\perp$、また $(1,2)$ の評価次元は上の完全交叉から16である。
従ってこの倍率空間は $(UV^2)^\perp=U$ であり、正確に元の第四項の像である。
第一ホモロジーも前の系で零なので、第二はEuler式により24となる。

正則なら $K$ は単純の解消そのもので、$I$ の一次関係はちょうどこの24個である。
さらに $n\ge7$ なら全項の係数が $D$ で全て完全である。
最小次数と、標準Hilbert関数の35に対する $S_4$ の次元20の差で四次生成は15個。
第三Torの唯一の可能次数6での次元はHilbert分子
$(1-z)^4(1+4z+10z^2+20z^3/(1-z))=1-15z^4+24z^5-10z^6$ で10となる。
第四Torは生成可能な次数が7以上なので存在しない。
\end{proof}
